\chapter{Fondations : réseaux, fonctions complètement monotones et vérification formelle}
\label{ch1:chap}

\begin{quote}\small\itshape
Ce chapitre pose le décor. Il ne suppose que le programme d'analyse et d'algèbre linéaire de première année de classe préparatoire,
plus les séries de Fourier de seconde année. Les démonstrations sont écrites à la main, comme on les rédigerait à un concours ;
la machine (Lean~4) n'intervient que dans la dernière partie, pour dire ce que signifie \og vérifié par ordinateur \fg{} et, surtout,
ce que cela ne signifie pas.
\end{quote}

\section*{Comment lire ce livre}

Chaque affirmation porte une étiquette de \emph{niveau} qui dit \emph{qui a vérifié quoi}.
Ces niveaux, empruntés à la méthode Elenchus~\cite{elenchus} et détaillés au §\ref{ch1:sec:tiers}, sont les suivants.
\begin{itemize}
\item \tierA{} : énoncé vérifié par le noyau de Lean~4, avec un contrôle des axiomes utilisés.
\item \tierL{} : résultat de la littérature, lu et référencé précisément.
\item \tierC{} : argument humain, écrit en détail dans ce livre, mais non vérifié par machine.
\item \tierX{} : exploration numérique en virgule flottante ; elle illustre, elle ne prouve jamais.
\end{itemize}

\begin{warnbox}[: honnêteté sur l'origine des résultats]
Tous les résultats de ce livre ont été produits par des agents d'IA (Claude) sous la direction de l'auteur, puis audités
par un \emph{autre} agent d'IA. \textbf{Aucun humain n'a audité ces résultats.} De plus, tout ce qui s'appuie sur le dépôt
\texttt{openai/math} est \emph{conditionnel} aux théorèmes de ce dépôt, qui sont eux-mêmes générés par IA et n'ont pas été
relus par des pairs~\cite{oai}. Un résultat de niveau \tierA{} garantit que l'énoncé \emph{Lean} découle des axiomes
autorisés ; il ne garantit pas que cet énoncé traduit fidèlement l'énoncé mathématique ou physique que l'on a en tête.
\end{warnbox}

\section{Réseaux du plan}
\label{ch1:sec:reseaux}

\subsection{Réseau, covolume, densité}

\begin{definition}[Réseau du plan]\label{ch1:def:reseau}
Soit $(u,v)$ une base de $\R^2$. Le \emph{réseau} engendré est $\Lambda=\Z u\oplus\Z v=\{mu+nv : (m,n)\in\Z^2\}$.
Son \emph{covolume} est $\operatorname{covol}(\Lambda)=|\det(u,v)|$, l'aire du parallélogramme
$P=\{su+tv : s,t\in[0,1[\}$, appelé \emph{cellule fondamentale}. Le covolume ne dépend pas de la base choisie
(deux bases diffèrent par une matrice de $\mathrm{GL}_2(\Z)$, de déterminant $\pm1$).
\end{definition}

On note $N_\Lambda(R)=\#\{x\in\Lambda : |x|\le R\}$ le nombre de points de $\Lambda$ dans le disque fermé de rayon $R$.
Un réseau est un exemple de \emph{configuration localement finie} : chaque disque contient un nombre fini de points.

\begin{proposition}[Densité d'un réseau]\label{ch1:prop:comptage}
Si $V=\operatorname{covol}(\Lambda)$, alors $N_\Lambda(R)/(\pi R^2)\to 1/V$ quand $R\to\infty$. Plus précisément, avec
$\delta=|u|+|v|$ et pour $R>\delta$,
\[
  \pi (R-\delta)^2\ \le\ V\,N_\Lambda(R)\ \le\ \pi (R+\delta)^2 .
\]
\tierC{}
\end{proposition}

\begin{humanproof}
Les translatés $x+P$, $x\in\Lambda$, \emph{pavent} le plan : tout point $y=su'+tv'$ s'écrit de façon unique
$y=x+p$ avec $x\in\Lambda$ (partie entière des coordonnées dans la base $(u,v)$) et $p\in P$. Chaque cellule
$x+P$ a pour aire $V$ et pour diamètre au plus $\delta$.
\begin{itemize}
\item Si $|x|\le R$, alors $x+P\subset \overline{D}(0,R+\delta)$. Les cellules sont disjointes, donc
$V\,N_\Lambda(R)\le \pi(R+\delta)^2$.
\item Si $|y|\le R-\delta$, le point $y$ appartient à une cellule $x+P$ avec $|y-x|\le\delta$, donc $|x|\le R$.
Ainsi $\overline{D}(0,R-\delta)\subset\bigcup_{|x|\le R}(x+P)$ et $\pi(R-\delta)^2\le V\,N_\Lambda(R)$.
\end{itemize}
On divise par $\pi R^2$ et on fait tendre $R$ vers l'infini.
\end{humanproof}

On dit que $\Lambda$ est de \emph{densité} $1/V$. Un réseau de covolume $1$ est de densité $1$ : c'est la normalisation
de tout le livre, qui permet de comparer des réseaux différents à nombre de points égal par unité d'aire.

\subsection{Le réseau triangulaire}

\begin{definition}[Réseau triangulaire $A$]\label{ch1:def:A}
Soit $a>0$ tel que $a^2=2/\sqrt3$. Le \emph{réseau triangulaire} (ou hexagonal) de covolume $1$ est
\[
  A=\Z\,u\oplus\Z\,v,\qquad u=a(1,0),\quad v=a\bigl(\tfrac12,\tfrac{\sqrt3}{2}\bigr).
\]
\end{definition}

Le choix de $a$ vient du covolume : $\det(u,v)=a^2\sqrt3/2=1$. Numériquement $a\approx 1{,}0746$.

\begin{proposition}[Normes dans $A$]\label{ch1:prop:normesA}
Pour $(m,n)\in\Z^2$ on a $|mu+nv|^2=a^2\,(m^2+mn+n^2)$. En particulier la distance minimale entre deux points de $A$ est $a$, atteinte
exactement pour $6$ vecteurs (le point central a donc $6$ plus proches voisins, aux sommets d'un hexagone régulier). \tierC{}
\end{proposition}

\begin{humanproof}
Comme $|u|^2=|v|^2=a^2$ et $u\cdot v=a^2/2$, on a $|mu+nv|^2=a^2(m^2+n^2+mn)$. Or
$m^2+mn+n^2=(m+\tfrac n2)^2+\tfrac34n^2\ge0$, avec égalité seulement pour $m=n=0$. Les solutions entières de
$m^2+mn+n^2=1$ vérifient $\tfrac34 n^2\le 1$, donc $n\in\{-1,0,1\}$ ; on résout chaque cas : $(\pm1,0)$, $(0,\pm1)$,
$\pm(1,-1)$. Il y en a six, et la valeur $1$ est la plus petite non nulle car $m^2+mn+n^2$ est un entier positif.
\end{humanproof}

Les premières valeurs de $k=m^2+mn+n^2$ et leurs multiplicités, comparées à celles du réseau carré $\Z^2$ (de covolume $1$ lui aussi),
sont données par le tableau~\ref{ch1:tab:couches}. Ces nombres ont été dénombrés par calcul sur machine \tierX{} ; ils sont classiques.

\begin{table}[ht]
\centering
\begin{tabular}{@{}llll@{}}
\toprule
Réseau & $k$ (valeur de la forme) & multiplicité & $|x|^2$ \\
\midrule
$A$ & $1,\ 3,\ 4,\ 7,\ 9$ & $6,\ 6,\ 6,\ 12,\ 6$ & $a^2k\approx 1{,}155,\ 3{,}464,\ 4{,}619,\ 8{,}083,\ 10{,}392$ \\
$\Z^2$ & $1,\ 2,\ 4,\ 5,\ 8$ & $4,\ 4,\ 4,\ 8,\ 4$ & $1,\ 2,\ 4,\ 5,\ 8$ \\
\bottomrule
\end{tabular}
\caption{Premières couches des réseaux $A$ et $\Z^2$ (covolume $1$).}
\label{ch1:tab:couches}
\end{table}

Le réseau triangulaire a des voisins plus éloignés ($|x|^2\approx1{,}155$ contre $1$) et plus nombreux ($6$ contre $4$) :
il \og équilibre \fg{} mieux la répulsion, et c'est ce que le théorème d'optimalité universelle va quantifier.

\subsection{Réseau dual}

\begin{definition}[Réseau dual]\label{ch1:def:dual}
Si $\Lambda=M\Z^2$ avec $M$ la matrice de colonnes $u,v$, le \emph{réseau dual} est
$\Lambda^*=\{y\in\R^2 : x\cdot y\in\Z\ \ \forall x\in\Lambda\}=M^{-T}\Z^2$. Il a pour covolume $1/\operatorname{covol}(\Lambda)$.
\end{definition}

\begin{lemma}[$A$ est semblable à son dual]\label{ch1:lem:autodual}
Le réseau $A^*$ est l'image de $A$ par une rotation ; en particulier $A$ et $A^*$ ont les mêmes normes, avec les mêmes multiplicités. \tierC{}
\end{lemma}

\begin{humanproof}
Notons $M=a\begin{pmatrix}1&1/2\\0&\sqrt3/2\end{pmatrix}$. Un calcul direct donne
$M^{-T}=\frac1a\begin{pmatrix}1&0\\-1/\sqrt3&2/\sqrt3\end{pmatrix}$. Ses deux colonnes $b_1,b_2$ vérifient
$|b_1|^2=|b_2|^2=\frac{1}{a^2}\cdot\frac43=\frac{\sqrt3}{2}\cdot\frac43=a^2$ et $b_1\cdot b_2=-\frac{1}{a^2}\cdot\frac23=-\frac{a^2}{2}$ :
deux vecteurs de norme $a$ faisant un angle de $120^\circ$. Alors $(b_1,b_1+b_2)$ est une base de $A^*$ formée de deux vecteurs
de norme $a$ faisant un angle de $60^\circ$, c'est-à-dire la base de $A$ tournée. $A^*$ est donc l'image de $A$ par une rotation.
\end{humanproof}

\subsection{Énergie d'un réseau}

On étudie des particules identiques qui se repoussent par un \emph{potentiel de paire} $g$, fonction du \emph{carré} $t=|x-y|^2$ de
la distance (convention de~\cite{ckmrv} et de~\cite{oai} : les gaussiennes $e^{-\alpha t}$ en sont alors les briques de base).
Le potentiel de Riesz $|x-y|^{-s}$ s'écrit $g(t)=t^{-s/2}$.

\begin{definition}[Énergie d'un réseau]\label{ch1:def:enerreseau}
Pour $g:\,]0,\infty[\to[0,\infty[$, l'énergie par particule d'un réseau $\Lambda$ est
$E_g(\Lambda)=\sum_{x\in\Lambda\setminus\{0\}} g(|x|^2)\in[0,+\infty]$.
\end{definition}

Pour une configuration localement finie $C$ de densité $1$, on définit l'\emph{énergie inférieure} par la moyenne sur les points
du disque de rayon $R$ des interactions avec les autres points du disque, puis une limite inférieure :
\[
  E_g(C;R)=\frac{1}{\#C_R}\sum_{x\in C_R}\ \sum_{y\in C_R\setminus\{x\}} g(|x-y|^2),\qquad E_g(C)=\liminf_{R\to\infty}E_g(C;R)\in[0,\infty],
\]
avec $C_R=C\cap\overline D(0,R)$. Chaque paire est comptée deux fois, ce qui donne bien $\sum_{a\neq0}g(|a|^2)$ pour un réseau.
\tierC{} La définition exacte utilisée par le dépôt \texttt{openai/math}, en arithmétique $[0,\infty]$, sera citée au §\ref{ch1:sec:univ}.

\begin{proposition}[Énergie de Riesz d'un réseau]\label{ch1:prop:riesz}
Pour $g(t)=t^{-s/2}$ avec $s>0$, l'énergie $E_g(\Lambda)=\sum_{x\ne0}|x|^{-s}$ est finie si et seulement si $s>2$. \tierC{}
\end{proposition}

\begin{humanproof}
Soit $d>0$ la plus petite norme non nulle de $\Lambda$ et $V$ son covolume. Les points non nuls de norme $\le d$ sont en nombre fini ;
pour $k\ge0$ notons $S_k=\{x\in\Lambda : 2^kd<|x|\le 2^{k+1}d\}$. Par la proposition~\ref{ch1:prop:comptage},
$\#S_k=N_\Lambda(2^{k+1}d)-N_\Lambda(2^kd)=\frac{3\pi d^2}{V}4^k+O(2^k)$, donc il existe $0<c<C$ avec $c\,4^k\le\#S_k\le C\,4^k$ pour $k\ge k_0$.
Sur $S_k$ on a $(2^{k+1}d)^{-s}\le |x|^{-s}<(2^kd)^{-s}$. La somme sur $S_k$ est donc comprise entre des multiples positifs de $4^k2^{-ks}=2^{k(2-s)}$.
La série $\sum_k 2^{k(2-s)}$ converge si et seulement si $s>2$.
\end{humanproof}

\begin{remark}
Pour $0<s\le2$ l'énergie d'un réseau est infinie : l'inégalité d'optimalité $\infty\le\infty$ est alors vraie mais sans contenu, et une comparaison
significative exigerait une \emph{régularisation}. \tierC{}
\end{remark}

\subsection{Séries thêta et formule de Poisson}
\label{ch1:sec:poisson}

Pour $\alpha>0$, la \emph{série thêta gaussienne} du réseau $\Lambda$ est
\[
  \Theta_\Lambda(\alpha)=\sum_{x\in\Lambda}e^{-\pi\alpha|x|^2}\ \ (>1).
\]
(Ne pas confondre avec les sommes de Gauss arithmétiques $\sum e^{2\pi i n^2/p}$ : il s'agit ici de sommes d'exponentielles
\emph{réelles} de formes quadratiques.) La série converge, par comparaison avec une série géométrique via la minoration
$|mu+nv|^2\ge c(m^2+n^2)$. L'énergie d'une interaction gaussienne $g(t)=e^{-\pi\alpha t}$ est donc
$E_g(\Lambda)=\Theta_\Lambda(\alpha)-1$.

Le résultat clé est que $\Theta_\Lambda$ se transforme simplement quand on passe de $\alpha$ à $1/\alpha$ et de $\Lambda$ à $\Lambda^*$.

\begin{theorem}[Formule de Poisson pour les gaussiennes]\label{ch1:thm:poisson}
Pour tout réseau $\Lambda$ du plan de covolume $V$ et tout $\alpha>0$,
\[
  \Theta_\Lambda(\alpha)=\frac{1}{V\,\alpha}\ \Theta_{\Lambda^*}\!\Bigl(\frac1\alpha\Bigr).
\]
Pour le réseau triangulaire ($V=1$, $A^*\simeq A$) : $\Theta_A(\alpha)=\alpha^{-1}\Theta_A(1/\alpha)$. \tierC{}\
(La convergence ponctuelle de la série de Fourier utilisée à l'étape~3 est le théorème de Fourier sur le tore : \tierL{}.)
\end{theorem}

\begin{humanproof}[Démonstration humaine en trois étapes]
Convention : $\hat f(\xi)=\int f(x)e^{-2\pi i\,\xi\cdot x}\,dx$.

\emph{Étape 1 : transformée de Fourier d'une gaussienne en dimension 1.} Soit $\varphi(x)=e^{-\pi\alpha x^2}$ et $h=\hat\varphi$.
Par dérivation sous l'intégrale (domination par $|x|e^{-\pi\alpha x^2}$, intégrable),
$h'(\xi)=\int(-2\pi i x)\varphi(x)e^{-2\pi i\xi x}dx$. Or $\varphi'(x)=-2\pi\alpha x\varphi(x)$, donc $-2\pi ix\varphi(x)=\frac{i}{\alpha}\varphi'(x)$ et
\[
  h'(\xi)=\frac i\alpha\int\varphi'(x)e^{-2\pi i\xi x}dx=\frac i\alpha\Bigl([\varphi e^{-2\pi i\xi x}]_{-\infty}^{+\infty}+2\pi i\xi\int\varphi(x)e^{-2\pi i\xi x}dx\Bigr)=-\frac{2\pi\xi}{\alpha}\,h(\xi),
\]
(le crochet est nul car $\varphi\to0$ aux bornes). Cette équation différentielle linéaire donne $h(\xi)=h(0)\,e^{-\pi\xi^2/\alpha}$, et
$h(0)=\int e^{-\pi\alpha x^2}dx=\alpha^{-1/2}$ (intégrale de Gauss $\int e^{-\pi x^2}dx=1$, calculée en coordonnées polaires). Ainsi
$\widehat{e^{-\pi\alpha x^2}}(\xi)=\alpha^{-1/2}e^{-\pi\xi^2/\alpha}$. En dimension $2$, par Fubini,
$\widehat{e^{-\pi\alpha|x|^2}}(\xi)=\alpha^{-1}e^{-\pi|\xi|^2/\alpha}$.

\emph{Étape 2 : périodisation.} Écrivons $\Lambda=M\Z^2$ et posons $f(x)=e^{-\pi\alpha|Mx|^2}$ pour $x\in\R^2$. La fonction
$F(x)=\sum_{n\in\Z^2}f(x+n)$ est bien définie (convergence normale locale, ainsi que celle de ses dérivées), $\mathcal C^\infty$, et $\Z^2$-périodique.
Ses coefficients de Fourier sont, pour $k\in\Z^2$,
\[
  c_k=\int_{[0,1]^2}F(x)e^{-2\pi ik\cdot x}dx=\int_{\R^2}f(x)e^{-2\pi ik\cdot x}dx=\hat f(k),
\]
(on réunit les translatés de $[0,1]^2$ pour obtenir $\R^2$). Le changement de variable $u=Mx$ ($dx=du/V$) donne
$\hat f(k)=\frac1V\int e^{-\pi\alpha|u|^2}e^{-2\pi i(M^{-T}k)\cdot u}du=\frac{1}{V\alpha}e^{-\pi|M^{-T}k|^2/\alpha}$ d'après l'étape~1.

\emph{Étape 3 : évaluation en $0$.} Les $c_k$ sont sommables sur $\Z^2$ et $F$ est continue ; la série de Fourier de $F$ converge donc normalement
vers une fonction continue de mêmes coefficients, qui est $F$ (unicité des coefficients, théorème de Fejér). En $x=0$ :
\[
  \sum_{n\in\Z^2}e^{-\pi\alpha|Mn|^2}=F(0)=\sum_{k\in\Z^2}c_k=\frac{1}{V\alpha}\sum_{k\in\Z^2}e^{-\pi|M^{-T}k|^2/\alpha}.
\]
Le membre de gauche est $\Theta_\Lambda(\alpha)$, et $M^{-T}\Z^2=\Lambda^*$ : le membre de droite est $\frac{1}{V\alpha}\Theta_{\Lambda^*}(1/\alpha)$.
Pour $A$, le lemme~\ref{ch1:lem:autodual} donne $\Theta_{A^*}=\Theta_A$.
\end{humanproof}

\begin{example}[Fonction thêta de Jacobi]\label{ch1:ex:jacobi}
En dimension~1, la même démonstration donne, pour $\vartheta(\alpha)=\sum_{n\in\Z}e^{-\pi\alpha n^2}$, la relation
$\vartheta(1/\alpha)=\sqrt\alpha\,\vartheta(\alpha)$. Elle montre que $\vartheta(\alpha)\sim\alpha^{-1/2}$ quand $\alpha\to0$ : pour un petit $\alpha$ la somme
\og voit \fg{} l'intégrale de Gauss. \tierC{}
\end{example}

\begin{numbox}[: la formule de Poisson avec ses contrôles (code \texttt{code/ch1\_theta.py})]
Pour $A$, on compare $\Theta_A(\alpha)$ à $\alpha^{-1}\Theta_A(1/\alpha)$ (contrôle positif), puis à $\Theta_A(1/\alpha)$ \emph{sans} le facteur
$\alpha^{-1}$ (contrôle négatif). Sortie réelle du script (en double précision) :
\begin{center}\small
\begin{tabular}{@{}rlll@{}}
\toprule
$\alpha$ & $\Theta_A(\alpha)$ & écart avec $\alpha^{-1}\Theta_A(1/\alpha)$ & écart sans le facteur \\
\midrule
$0{,}5$ & $2{,}008477918598$ & $4{,}4\cdot10^{-16}$ & $1{,}0$ \\
$2$ & $1{,}004238959299$ & $2{,}2\cdot10^{-16}$ & $1{,}0$ \\
$3$ & $1{,}000112671200$ & $8{,}9\cdot10^{-16}$ & $2{,}0$ \\
\bottomrule
\end{tabular}
\end{center}
Le contrôle négatif est discriminant ($\alpha\neq1$) ; en $\alpha=1$ les deux écarts sont nuls, donc ce point ne teste rien.
Le script compare aussi $\Theta_A(\alpha)\le\Theta_{\Z^2}(\alpha)$ sur huit valeurs de $\alpha\in[0{,}1;5]$ (toujours vrai), par exemple
$\Theta_A(1)=1{,}159595267$ contre $\Theta_{\Z^2}(1)=1{,}180340599$. \tierX{} : cela illustre le théorème de Montgomery ci-dessous, sans le prouver.
\end{numbox}

\begin{theorem}[Montgomery, 1988]\label{ch1:thm:montgomery}
Pour tout réseau $\Lambda$ du plan de covolume $1$ et tout $\alpha>0$, $\Theta_A(\alpha)\le\Theta_\Lambda(\alpha)$, avec égalité si et seulement si $\Lambda$ est
l'image de $A$ par une rotation. \tierL{}~\cite{montgomery}
\end{theorem}

Montgomery résout ainsi le problème d'optimalité \emph{parmi les réseaux}, pour toute interaction gaussienne. Nous y reviendrons au §\ref{ch1:sec:univ}
pour passer des gaussiennes aux potentiels complètement monotones. La difficulté est ailleurs : comparer $A$ non pas à d'autres réseaux, mais à
\emph{toutes} les configurations de points de densité $1$.

\section{Fonctions complètement monotones}
\label{ch1:sec:cm}

\subsection{Définition et premiers exemples}

\begin{definition}[Fonction complètement monotone]\label{ch1:def:cm}
Une fonction $g:\,]0,\infty[\to\R$ est \emph{complètement monotone} (CM) si elle est de classe $\mathcal C^\infty$ et si
$(-1)^n g^{(n)}(t)\ge0$ pour tout $n\in\N$ et tout $t>0$.
\end{definition}

Le cas $n=0$ dit $g\ge0$ ; le cas $n=1$ dit que $g$ est décroissante ; le cas $n=2$ dit que $g$ est convexe ; et ainsi de suite.
Les inégalités sont \emph{larges} : les constantes sont CM. Dans le dépôt \texttt{openai/math}, une fonction qui est $\mathcal C^\infty$ sur $]0,\infty[$,
positive et CM est appelée \emph{admissible}~\cite{oai} ; comme $n=0$ donne la positivité, les deux notions coïncident. Cette coïncidence est
elle-même un théorème Lean : \lean{ProgCM.isCompletelyMonotoneOn_iff_admissible} \tierA{} (bibliothèque du programme P3~\cite{p3}).

\begin{example}[Exponentielle, puissances]\label{ch1:ex:exemples}
\begin{enumerate}
\item $g(t)=e^{-at}$, $a\ge0$ : $(-1)^ng^{(n)}(t)=a^ne^{-at}\ge0$. \tierC{}
\item $g(t)=t^{-s}$, $s\ge0$ : $(-1)^ng^{(n)}(t)=s(s+1)\cdots(s+n-1)\,t^{-s-n}\ge0$ (récurrence immédiate). \tierC{}
Pour $s=\frac12$ c'est la répulsion de Coulomb $1/r$ en la variable $t=r^2$. La CM de $t\mapsto t^{-s/2}$ (Riesz) est aussi
vérifiée par le noyau :\lean{TriangularRiesz.riesz_admissible} \tierA{}.
\end{enumerate}
\end{example}

\begin{example}[Contre-exemples]\label{ch1:ex:contre}
$g(t)=t$ n'est pas CM ($-g'=-1<0$), ni $\cos t$ (elle prend des valeurs négatives), ni $e^{t}$ ($-g'<0$).
Un contre-exemple plus subtil, d'origine physique (combinaison charge--quadrupôle) est
$q(t)=t^{-1/2}-2(t+1)^{-1/2}+(t+4)^{-1/2}$ : on a $q(3)=3^{-1/2}-2\cdot4^{-1/2}+7^{-1/2}\approx0{,}5774-1+0{,}3780<0$,
donc $q$ n'est pas positive, a fortiori pas CM. La vérification rigoureuse de $q(3)<0$ et la non-CM de $q$ sont vérifiées par le noyau
(\lean{ProgCM.quadrupole_not_cm}) \tierA{}.
\end{example}

\subsection{Propriétés de clôture : preuves humaines}

\begin{proposition}[Somme, produit]\label{ch1:prop:cloture}
Soient $f,g$ CM et $c\ge0$. Alors $f+g$, $cf$ et $fg$ sont CM. \tierC{} (version Lean : le programme P3~\cite{p3}, \tierA{})
\end{proposition}

\begin{humanproof}
La somme et la multiplication par une constante positive sont immédiates par linéarité de la dérivation.
Pour le produit, la formule de Leibniz donne
\[
  (-1)^n(fg)^{(n)}=\sum_{k=0}^{n}\binom nk\ \bigl[(-1)^kf^{(k)}\bigr]\ \bigl[(-1)^{n-k}g^{(n-k)}\bigr],
\]
car $(-1)^n=(-1)^k(-1)^{n-k}$. Chaque crochet est positif, donc la somme aussi.
\end{humanproof}

Une \emph{fonction de Bernstein} est la bonne notion pour composer avec une fonction CM.

\begin{definition}[Fonction de Bernstein]\label{ch1:def:bernstein}
Une fonction $h:\,]0,\infty[\to\,]0,\infty[$ de classe $\mathcal C^\infty$ est \emph{de Bernstein} si sa dérivée $h'$ est complètement monotone.
(En particulier $h$ est croissante.)
\end{definition}

\begin{example}\label{ch1:ex:bernstein}
Sont de Bernstein : $h(t)=at+b$ ($a\ge0$, $b>0$ ou $a>0$, $b\ge 0$) ; $h(t)=t^\alpha$ pour $0<\alpha\le1$, car
$h'(t)=\alpha\,t^{-(1-\alpha)}$ est CM par l'exemple~\ref{ch1:ex:exemples} ; $h(t)=1-e^{-at}$ ($a>0$), car $h'=ae^{-at}$ ; $h(t)=\ln(1+t)$,
car $h'(t)=1/(1+t)$. En particulier $h(t)=\sqrt t$ (cas $\alpha=\frac12$) l'est. \tierC{}
\end{example}

\begin{theorem}[Subordination : CM $\circ$ Bernstein]\label{ch1:thm:comp}
Si $f$ est CM et $h$ est de Bernstein, alors $f\circ h$ est CM. \tierC{} Version Lean :
\lean{ProgCM.IsCompletelyMonotoneOn.comp_bernstein} \tierA{}, dont la démonstration suit exactement la récurrence ci-dessous~\cite{p3}.
\end{theorem}

\begin{humanproof}
La composée $f\circ h$ est bien définie (car $h>0$) et de classe $\mathcal C^\infty$. On montre par récurrence forte sur $n$ la propriété
\[
  \mathcal P(n) :\quad \text{pour toute } f \text{ CM et toute } h \text{ de Bernstein},\ \ (-1)^n(f\circ h)^{(n)}\ge0\ \text{sur }]0,\infty[.
\]
Il est essentiel de quantifier sur \emph{toutes} les fonctions $f$ et $h$ : l'hypothèse de récurrence sera appliquée à une autre fonction que $f$.

\emph{Initialisation.} $\mathcal P(0)$ : $f\circ h\ge0$ car $f\ge0$.

\emph{Hérédité.} Supposons $\mathcal P(k)$ pour tout $k\le n$. Soit $f$ CM et $h$ de Bernstein. On a $(f\circ h)'=(f'\circ h)\,h'$, donc
\[
  -(f\circ h)'=G\cdot h',\qquad G=(-f')\circ h .
\]
Or $-f'$ est CM (c'est $f$ dérivée une fois, avec le signe : $(-1)^m(-f')^{(m)}=(-1)^{m+1}f^{(m+1)}\ge0$). Par $\mathcal P(k)$ appliquée à la fonction CM $-f'$,
$(-1)^kG^{(k)}\ge0$ pour tout $k\le n$. D'autre part $h'$ est CM par définition. En dérivant $n$ fois et par la formule de Leibniz,
\[
  (-1)^{n+1}(f\circ h)^{(n+1)}=(-1)^{n}\bigl(G\,h'\bigr)^{(n)}=\sum_{k=0}^{n}\binom nk\ \bigl[(-1)^kG^{(k)}\bigr]\ \bigl[(-1)^{n-k}(h')^{(n-k)}\bigr]\ \ge0 .
\]
D'où $\mathcal P(n+1)$.
\end{humanproof}

\begin{corollary}[Le cas de la racine carrée]\label{ch1:cor:sqrt}
Si $f$ est CM, alors $t\mapsto f(\sqrt t)$ est CM. En particulier, pour $\kappa\ge0$, $t\mapsto e^{-\kappa\sqrt t}$ est CM, et il en est de même du
potentiel de Yukawa planaire $e^{-\kappa r}/r$ écrit en la variable $t=r^2$ : $g_\kappa(t)=e^{-\kappa\sqrt t}\,t^{-1/2}$. \tierC{}
(Lean : \lean{ProgCM.IsCompletelyMonotoneOn.comp_sqrt}, \lean{ProgCM.phi_cm} \tierA{}.)
\end{corollary}

\begin{humanproof}
Le théorème~\ref{ch1:thm:comp} avec $h=\sqrt{\ \cdot\ }$ (exemple~\ref{ch1:ex:bernstein}) donne la première phrase. Avec $f(u)=e^{-\kappa u}$ on obtient que
$e^{-\kappa\sqrt t}$ est CM. Le potentiel de Yukawa est le \emph{produit} de cette fonction par $t^{-1/2}$, CM d'après l'exemple~\ref{ch1:ex:exemples} ;
la proposition~\ref{ch1:prop:cloture} conclut.
\end{humanproof}

\begin{remark}[Vérification directe sur un cas]
Pour $\psi(t)=e^{-\sqrt t}$ on a $\psi'(t)=-\frac{1}{2\sqrt t}e^{-\sqrt t}$ et $\psi''(t)=\bigl(\frac1{4t}+\frac1{4t^{3/2}}\bigr)e^{-\sqrt t}$. Les signes
$-\psi'\ge0$ et $\psi''\ge0$ sont ceux annoncés. Une machine vérifie les signes jusqu'à l'ordre $8$ en sept points
(script \texttt{code/ch1\_theta.py}) : $e^{-\sqrt t}$ passe, $\cos t$ (contrôle négatif) échoue à chaque ordre \tierX{}. Ce n'est pas une preuve :
la preuve est celle de la récurrence ci-dessus.
\end{remark}

\begin{physbox}[: pourquoi les physiciens rencontrent des fonctions CM]
Un potentiel qui décroît avec la distance en s'amortissant \og de plus en plus lentement \fg{} (Coulomb, Yukawa écranté, dipôle) est typiquement CM.
Les cinq potentiels physiques du tableau~\ref{ch1:tab:flux} comprennent les excitons dipolaires $2\bigl(1/r-1/\sqrt{r^2+d^2}\bigr)$ (doi:\,10.5281/zenodo.23646103) et le potentiel de Yukawa
(doi:\,10.5281/zenodo.23646105). Dans les deux cas, l'admissibilité est un théorème de niveau \tierA{} ; la conséquence physique (le réseau triangulaire est
d'énergie minimale) reste conditionnelle à \texttt{openai/math}.
\end{physbox}

\subsection{Le théorème de Bernstein et les mélanges de gaussiennes}

\begin{theorem}[Bernstein--Widder]\label{ch1:thm:bernstein}
Une fonction $g:\,]0,\infty[\to\R$ est complètement monotone si et seulement s'il existe une mesure positive $\mu$ sur $[0,\infty[$ telle que
\[
  g(t)=\int_{[0,\infty[}e^{-\alpha t}\,d\mu(\alpha)\qquad(t>0).
\]
La mesure $\mu$ est alors unique. \tierL{}~\cite{bernstein,widder}
\end{theorem}

La lecture géométrique est essentielle : avec $t=|x|^2$, on a $g(|x|^2)=\int e^{-\alpha|x|^2}d\mu(\alpha)$, c'est-à-dire que \emph{tout potentiel complètement monotone
est un mélange positif de gaussiennes}. Le sens \og réciproque \fg{} est facile et nous le démontrons ; le sens direct (CM $\Rightarrow$ mélange) est un résultat
profond, que nous citons, et qui n'est \emph{pas} formalisé dans le travail décrit ici.

\begin{proposition}[Les mélanges sont CM (sens facile)]\label{ch1:prop:facile}
Soit $\mu$ une mesure positive sur $[0,\infty[$ telle que $\int e^{-\alpha t}d\mu(\alpha)<\infty$ pour tout $t>0$. Alors $g(t)=\int e^{-\alpha t}d\mu(\alpha)$ est CM.
\tierC{} (Lean : \lean{ProgCM.laplace_measure} \tierA{}, énoncé de ce même sens.)
\end{proposition}

\begin{humanproof}
Fixons $t_0>0$. Pour $t\ge t_0$ et $n\in\N$, on a $0\le\alpha^ne^{-\alpha t}\le C_n\,e^{-\alpha t_0/2}$ avec
$C_n=\sup_{\alpha\ge0}\alpha^ne^{-\alpha t_0/2}<\infty$, et $\alpha\mapsto e^{-\alpha t_0/2}$ est $\mu$-intégrable par hypothèse. Le théorème de dérivation sous le signe
somme (domination uniforme sur $[t_0,\infty[$) montre, par récurrence sur $n$, que $g$ est de classe $\mathcal C^\infty$ et
$(-1)^ng^{(n)}(t)=\int\alpha^ne^{-\alpha t}d\mu(\alpha)\ge0$.
\end{humanproof}

\begin{example}[Un potentiel CM que les règles de clôture ne voient pas]\label{ch1:ex:dipole}
Pour $c>0$, $\varphi(t)=t^{-1/2}-(t+c)^{-1/2}$ est une \emph{différence} de deux fonctions CM, et rien dans la proposition~\ref{ch1:prop:cloture} ne
permet de conclure. Mais la formule d'Euler $u^{-1/2}=\frac{1}{\sqrt\pi}\int_0^\infty \alpha^{-1/2}e^{-\alpha u}d\alpha$ ($u>0$) donne
\[
  \varphi(t)=\frac{1}{\sqrt\pi}\int_0^\infty\alpha^{-1/2}\bigl(1-e^{-c\alpha}\bigr)\,e^{-\alpha t}\,d\alpha ,
\]
un mélange à poids \emph{positif} ($1-e^{-c\alpha}>0$) : $\varphi$ est CM par la proposition~\ref{ch1:prop:facile}. C'est le potentiel de l'exciton
dipolaire, au facteur $2$ près. \tierC{}
\end{example}

\paragraph{Passage de l'énergie des gaussiennes à celle des mélanges.} Soit $g(t)=\int e^{-\pi\alpha t}d\nu(\alpha)$ un mélange, $\nu$ positive.
Pour une configuration finie $C_R$, par positivité (théorème de Tonelli) la double somme se met sous l'intégrale :
\[
  E_g(C;R)=\int E_{e^{-\pi\alpha\,\cdot}}(C;R)\,d\nu(\alpha).
\]
Le lemme de Fatou (valable dans $[0,\infty]$) donne alors $E_g(C)\ge\int\liminf_RE_{e^{-\pi\alpha\cdot}}(C;R)\,d\nu(\alpha)=\int E_{e^{-\pi\alpha\cdot}}(C)\,d\nu(\alpha)$.
Si on sait que, pour chaque $\alpha$, aucune configuration n'a une énergie gaussienne inférieure à celle de $A$, on en déduit
$E_g(C)\ge\int(\Theta_A(\alpha)-1)\,d\nu(\alpha)=E_g(A)$, par Tonelli sur la somme du réseau. \tierC{}
\emph{Le cas gaussien est donc le c\oe ur du problème} ; c'est lui, pour toute valeur de $\alpha$, que le théorème de la partie suivante affirme.
La même idée, appliquée aux \emph{réseaux} avec le théorème de Montgomery~\ref{ch1:thm:montgomery}, donne : $E_g(A)\le E_g(\Lambda)$ pour tout réseau $\Lambda$ de
covolume $1$ et tout $g$ CM (mélange $\nu$ de gaussiennes), car $E_g(\Lambda)=\int(\Theta_\Lambda(\alpha)-1)\,d\nu(\alpha)$. \tierC{}\ (conditionnelle à~\ref{ch1:thm:montgomery} et à~\ref{ch1:thm:bernstein}, \tierL{})

\section{L'optimalité universelle}
\label{ch1:sec:univ}

\subsection{Énoncé}

Cohn et Kumar~\cite{ck} ont défini une notion d'\emph{optimalité universelle} : une configuration est universellement optimale si elle minimise l'énergie pour
\emph{tout} potentiel complètement monotone de la distance au carré, simultanément. Ils ont conjecturé que cela se produit pour de nombreuses configurations
remarquables, dont le réseau triangulaire dans le plan. Cohn, Kumar, Miller, Radchenko et Viazovska~\cite{ckmrv} ont démontré la conjecture pour le réseau $E_8$ en dimension $8$ et le réseau de Leech
en dimension $24$ \tierL{}, en s'appuyant sur le travail de Viazovska~\cite{viazovska}. La dimension $2$ restait ouverte dans la littérature que nous avons lue.

Le dépôt \texttt{openai/math} contient un énoncé et une démonstration en Lean~4 du cas planaire~\cite{oai}. Voici l'énoncé tel qu'il figure dans le fichier de
défi, en français puis en Lean.

\begin{theorem}[Optimalité universelle du réseau triangulaire (\texttt{openai/math})]\label{ch1:thm:univ}
Soit $g$ admissible (définition~\ref{ch1:def:cm}) et $C\subset\R^2$ une configuration localement finie de densité centrée $1$. Alors
\[
  \mathcal E_g(A)\ \le\ E_g(C)\qquad\text{et}\qquad \mathcal E_g(A)=E_g(A),\qquad\text{où}\quad \mathcal E_g(A)=\sum_{x\in A\setminus 0}g(|x|^2),
\]
avec des valeurs dans $[0,\infty]$, et $E_g(\cdot)$ l'énergie inférieure sur les disques fermés centrés (Lean : \lean{latticeEnergy g} pour $\mathcal E_g(A)$ et \lean{energy g} pour $E_g$). \tierA{} \emph{pour l'énoncé Lean}
\lean{OAI.AtomicTriangular.universal_energy_minimum}, \textbf{conditionnel} : la démonstration est générée par IA et n'a pas été relue par des pairs.
\end{theorem}

\begin{leanbox}[: l'énoncé principal (extrait du fichier de défi, la preuve y est \lean{sorry} par construction)]
\begin{lstlisting}[language=Lean]
theorem universal_energy_minimum (g : ℝ → ℝ) (C : Set Plane)
    (hg : AdmissiblePotential g) (hC : LocallyFinite C) (hd : DensityOne C) :
    latticeEnergy g ≤ energy g C ∧ latticeEnergy g = energy g A
\end{lstlisting}
\end{leanbox}

Trois remarques de lecture, essentielles pour la suite.
\begin{enumerate}
\item L'énergie est une \emph{limite inférieure} sur des disques \emph{fermés centrés en $0$}. La définition de CKMRV~\cite{ckmrv} exige en plus que la limite existe pour la
configuration optimale. La notion utilisée ici est donc plus faible \tierC{}.
\item Le théorème ne dit pas que $A$ est le seul minimiseur, et l'unicité est fausse en général : le corollaire D3 du livre v1~\cite{v1} établit la
non-unicité en retirant un point de $A$ (une \og lacune \fg{}), ce qui ne change ni la densité ni l'énergie inférieure. \tierA{} (conditionnel)
\item Pour $s\le2$ (Riesz), l'énergie de $A$ est infinie (proposition~\ref{ch1:prop:riesz}) : l'inégalité est alors vraie sans contenu. Le contenu apparaît pour $s>2$.
\end{enumerate}

\subsection{Le fil conducteur des formalisations}

La suite du livre suit un enchaînement naturel : chaque étape prend le résultat précédent comme une boîte noire.
Le tableau~\ref{ch1:tab:flux} le résume. Les numéros de DOI et les niveaux sont ceux des documents d'enregistrement.

\begin{table}[ht]
\centering\small
\begin{tabular}{@{}p{0.04\textwidth}p{0.50\textwidth}p{0.17\textwidth}p{0.20\textwidth}@{}}
\toprule
\# & Étape & Niveau & Référence \\
\midrule
1 & Théorème amont (\texttt{openai/math}), cas planaire & \tierL{} et \tierA{} conditionnel & \cite{oai} \\
2 & Corollaires D1 à D4 : Riesz, toute densité, sens CKMRV et non-unicité, disques ouverts ou fermés & \tierA{}, conditionnel à 1 & \cite{v1} (livre v1) \\
3 & Cinq potentiels physiques (excitons dipolaires, Yukawa, Rytova--Keldysh, enstrophie $\sigma=3$, structure de Bogoliubov) & \tierA{} ; Keldysh$\leftrightarrow$Struve : \tierC{} & 5 dépôts Zenodo \\
4 & Bibliothèque CM/Bernstein (P3) : subordination, équivalence avec \lean{AdmissiblePotential} & \tierA{} & \cite{p3} \\
5 & Surprise quantique : l'énergie de point zéro, QU-1 faux pour 8 potentiels CM sur 12 ; certificat à $a=1$ ; pas de transition triangulaire$\to$oblique & \tierX{}, puis \tierA{} (logique) et \tierC{} (intégrales en boules) ; résultat négatif & Zenodo \\
6 & Sandier--Serfaty (réseaux d'Abrikosov) : audit, portage partiel en Lean, échec de l'extension au cas~3D & \tierC{}, \tierA{} partiel ; résultats négatifs & Zenodo \\
7 & Navier--Stokes (corpus \texttt{openai/NavierStokesAndEuler}) : reconstruction et rejeu noyau & \tierA{} ; pas de \og problème Clay résolu \fg{} & article en préparation \\
8 & Génétique des populations : fixation de Moran, équation de Price, théorème isotherme & \tierA{} & Zenodo \\
\bottomrule
\end{tabular}
\caption{L'enchaînement des formalisations, des fondations (ce chapitre) aux applications.}
\label{ch1:tab:flux}
\end{table}

\begin{warnbox}[: ce que le tableau n'affirme pas]
Aucune ligne ne signifie \og conjecture démontrée \fg{} au-delà de son niveau. Les lignes 5 et 6 contiennent des résultats \emph{négatifs} : elles
disent que certaines extensions plausibles échouent (ou ne sont pas établies), et c'est une information de même valeur que les résultats positifs.
L'étape~7 est un audit d'un corpus de tiers : le mérite revient à ses auteurs. Nous ne prétendons pas qu'un problème du Clay est résolu.
\end{warnbox}

\section{Lean 4 pour l'étudiant}
\label{ch1:sec:lean}

\subsection{Propositions et preuves}

Lean~4~\cite{lean4} est à la fois un langage de programmation et un \emph{assistant de preuve}. Pour un étudiant de classe préparatoire,
le dictionnaire suivant suffit pour lire les énoncés de ce livre (c'est l'idée de Curry et Howard).
\begin{itemize}
\item Une \emph{proposition} est un énoncé mathématique ; en Lean, c'est un objet de type \lean{Prop}.
\item Une \emph{preuve} de la proposition $P$ est un objet \lean{h : P}. Le noyau de Lean vérifie qu'un terme a bien le type annoncé.
\item $P\wedge Q$ s'écrit \lean{P ∧ Q} : une preuve est une \emph{paire} $\langle$\lean{hP}, \lean{hQ}$\rangle$.
\item $P\Rightarrow Q$ s'écrit \lean{P → Q} : une preuve est une \emph{fonction} qui transforme toute preuve de $P$ en une preuve de $Q$.
\item $\forall n,\ P(n)$ s'écrit \lean{∀ n, P n} : une preuve est une fonction qui, à chaque $n$, associe une preuve de $P(n)$.
\end{itemize}
La bibliothèque Mathlib~\cite{mathlib} fournit de nombreux théorèmes déjà démontrés (analyse, algèbre).

\subsection{Un exemple minimal, avec sa sortie réelle}

Le fichier suivant a été compilé tel quel avec Lean 4.34.0-rc2, sans aucune bibliothèque. Il contient trois théorèmes vrais et deux \og théorèmes \fg{}
faux $1=2$ que l'on ne démontre qu'en trichant : l'un avec \lean{sorry}, l'autre avec un axiome ajouté par l'utilisateur.

\begin{leanbox}[: \texttt{T.lean} (compilé le 2026-10-26)]
\begin{lstlisting}[language=Lean]
theorem and_swap' (p q : Prop) (h : p ∧ q) : q ∧ p := ⟨h.2, h.1⟩

theorem add_comm' (a b : Nat) : a + b = b + a := Nat.add_comm a b

theorem two_mul' (n : Nat) : 2 * n = n + n := by
  omega

theorem bad : 1 = 2 := by sorry

axiom cheat : False
theorem bad2 : 1 = 2 := cheat.elim

#print axioms and_swap'
#print axioms two_mul'
#print axioms bad
#print axioms bad2
\end{lstlisting}
\end{leanbox}

Les deux premiers théorèmes sont des \emph{termes de preuve} : on écrit directement l'objet. Le troisième utilise une \emph{tactique} (\lean{by omega}) :
un petit programme qui \emph{construit} le terme de preuve ; le noyau vérifie ensuite le terme produit. Voici la sortie réelle :

\begin{lstlisting}[language={}]
T.lean:8:8: warning: declaration uses `sorry`
'and_swap'' does not depend on any axioms
'two_mul'' depends on axioms: [propext, Quot.sound]
'bad' depends on axioms: [sorryAx]
'bad2' depends on axioms: [cheat]
\end{lstlisting}

Le code de retour du compilateur était $0$ : \textbf{Lean accepte un fichier qui contient \lean{sorry}}, en émettant seulement un avertissement.
\emph{Un code de retour nul ne signifie pas \og tout est démontré \fg{}.}

\subsection{Les axiomes, et la porte d'entrée}

Toute preuve en mathématiques repose sur des axiomes. Les mathématiques classiques de Mathlib en utilisent trois : \lean{propext} (deux propositions
équivalentes sont égales), \lean{Classical.choice} (axiome du choix) et \lean{Quot.sound} (compatibilité des quotients). Nous les appelons la \emph{liste blanche} :
\[
  \{\texttt{propext},\ \texttt{Classical.choice},\ \texttt{Quot.sound}\}.
\]
Tout autre axiome dans l'empreinte d'un théorème le disqualifie. Dans la sortie ci-dessus, l'empreinte de \lean{bad} contient \lean{sorryAx} (l'axiome interne qui
se cache derrière \lean{sorry}) et celle de \lean{bad2} contient l'axiome utilisateur \lean{cheat} : les deux faux théorèmes sont donc \emph{détectés}.
C'est le \emph{contrôle négatif} de la porte d'axiomes ; le \emph{contrôle positif} est que \lean{and_swap'} et \lean{two_mul'} passent. Une mesure sans ces deux contrôles ne prouve rien.

\begin{warnbox}[: \texttt{\#print axioms} ne suffit pas]
L'audit A2~\cite{a2} du corpus \texttt{openai/NavierStokesAndEuler} (2\,440 modules reconstruits, rejeu noyau de $53\,752+34\,648$ déclarations) a mis en évidence qu'une déclaration
\emph{ajoutée directement au noyau en désactivant sa vérification de types} (option \lean{debug.skipKernelTC}) est affichée par \lean{#print axioms} comme
\og ne dépend d'aucun axiome \fg{}, alors qu'elle peut établir \lean{False}. Le contrôle positif de l'audit (une telle fausse déclaration, fabriquée exprès) a été
\emph{invisible} à \lean{#print axioms} et \emph{attrapé} par le rejeu complet du noyau (\lean{leanchecker}) \tierA{}. La leçon : l'audit d'axiomes est nécessaire,
non suffisant.
\end{warnbox}

La porte de niveau \tierA{} de ce livre comprend donc \emph{quatre} contrôles :
\begin{enumerate}
\item compilation du projet sans erreur ;
\item liste d'axiomes de chaque théorème retenu incluse dans la liste blanche ;
\item \emph{rejeu noyau} complet des déclarations par l'outil \lean{leanchecker}, qui revérifie chaque terme sans faire confiance au compilateur ;
\item recherche textuelle des mots \lean{skipKernelTC}, \lean{addDecl}, \lean{native_decide} et \lean{implemented_by} dans les sources, qui sont les portes dérobées connues.
\end{enumerate}
S'y ajoute une règle d'organisation : celui qui \emph{produit} une preuve n'est pas celui qui la \emph{vérifie} (ici, deux agents d'IA distincts, et jamais un humain).

\subsection{Lire l'énoncé amont ; un corollaire court}

Une fois le théorème principal disponible, un corollaire se réduit souvent à vérifier une hypothèse. Voici le cas de Riesz,
cité depuis le fichier \lean{formal_cert/Triangular/Riesz.lean} du dépôt (les points de suspension marquent une partie omise) :

\begin{leanbox}[: corollaire de Riesz (extrait)]
\begin{lstlisting}[language=Lean]
noncomputable def riesz (s : ℝ) : ℝ → ℝ := fun t => t ^ (-(s / 2))

theorem riesz_admissible {s : ℝ} (hs : 0 < s) : AdmissiblePotential (riesz s) := by
  ...

theorem triangular_riesz_optimal {s : ℝ} (hs : 0 < s) (C : Set Plane)
    (hC : LocallyFinite C) (hd : DensityOne C) :
    latticeEnergy (riesz s) ≤ energy (riesz s) C :=
  (universal_energy_minimum (riesz s) C (riesz_admissible hs) hC hd).1
\end{lstlisting}
\end{leanbox}

La dernière ligne \og applique \fg{} le théorème principal à $g=\text{riesz}\ s$ en fournissant les trois preuves d'hypothèses, puis prend la première composante
\lean{.1} de la conjonction. C'est exactement une \emph{composition de fonctions} au sens du dictionnaire de Curry--Howard.
L'hypothèse \lean{riesz_admissible} est, elle, la propriété de complète monotonie de l'exemple~\ref{ch1:ex:exemples} (2).

\subsection{La récurrence du §\ref{ch1:sec:cm}, côté Lean}

Le théorème de subordination~\ref{ch1:thm:comp} a été formalisé~\cite{p3}. L'énoncé général et son cas particulier de la racine carrée sont les suivants
(extraits du fichier \lean{formal_cert/Programmes/CompletelyMonotone/Bernstein.lean}, preuves omises) :

\begin{leanbox}[: subordination (énoncés)]
\begin{lstlisting}[language=Lean]
theorem IsCompletelyMonotoneOn.comp_bernstein {f g : ℝ → ℝ} (hf : IsCompletelyMonotoneOn f)
    (hg : IsBernsteinOn g) (hpos : ∀ t, 0 < t → 0 < g t) :
    IsCompletelyMonotoneOn (fun t => f (g t))

theorem IsCompletelyMonotoneOn.comp_sqrt {f : ℝ → ℝ} (hf : IsCompletelyMonotoneOn f) :
    IsCompletelyMonotoneOn (fun t => f (Real.sqrt t))
\end{lstlisting}
\end{leanbox}

On y reconnaît la structure du théorème humain. Les hypothèses \lean{hf}, \lean{hg} et \lean{hpos} sont les trois
hypothèses du théorème ($f$ CM, $h$ de Bernstein, $h>0$), \lean{fun t => f (g t)} est la fonction $f\circ h$, et \lean{comp_sqrt} s'obtient du cas général en
fournissant la preuve que $\sqrt{\ \cdot\ }$ est de Bernstein et strictement positive sur $]0,\infty[$. Le dépôt indique que la preuve procède par récurrence
forte sur l'ordre, quantifiée sur toutes les fonctions CM $f$, avec la formule $(f\circ h)'=-((-f')\circ h)\cdot h'$ et le lemme de Leibniz : c'est
le même raisonnement que la démonstration humaine. \tierA{} pour les énoncés ci-dessus, \tierC{} pour notre rédaction humaine, dont la correspondance avec la
preuve Lean est un commentaire de l'auteur, non un théorème.

\begin{remark}[Pourquoi \lean{hpos} apparaît en Lean]
En mathématiques de papier, on sous-entend que $f\circ h$ a un sens. Lean ne le sous-entend pas : $f$ est une fonction définie partout sur $\R$, mais
n'est contrainte que sur $]0,\infty[$. Si $h\equiv0$ (qui est de Bernstein), $f\circ h$ serait la constante parasite $f(0)$, sur laquelle l'hypothèse ne dit rien. L'hypothèse
$h>0$ n'est pas du pinaillage : elle est ce qui rend l'énoncé \emph{vrai}.
\end{remark}

\begin{warnbox}[: récapitulatif des niveaux de ce chapitre]
\begin{itemize}
\item \tierL{} : formule de Poisson sur le tore (théorème de Fourier), théorème de Montgomery, théorème de Bernstein--Widder (sens direct), résultat de CKMRV en dimensions $8$ et $24$.
\item \tierC{} : toutes les démonstrations humaines de ce chapitre (comptage, énergie de Riesz, formule de Poisson, clôture des fonctions CM, subordination, mélanges, réduction de Fatou).
\item \tierA{} : les énoncés Lean cités (équivalence avec \lean{AdmissiblePotential}, subordination, \lean{laplace_measure}, \lean{quadrupole_not_cm}, l'énoncé amont), le dernier étant \emph{conditionnel}.
\item \tierX{} : le script \texttt{code/ch1\_theta.py}, qui illustre sans prouver.
\end{itemize}
\end{warnbox}

\subsection{Les niveaux Elenchus}
\label{ch1:sec:tiers}

La méthode Elenchus~\cite{elenchus} attache à chaque énoncé du livre l'un des quatre niveaux, selon la \emph{nature de la preuve disponible}.

\begin{center}\small
\begin{tabular}{@{}p{0.09\textwidth}p{0.40\textwidth}p{0.38\textwidth}@{}}
\toprule
Niveau & Ce qui est garanti & Ce qui ne l'est pas \\
\midrule
\tierA{} & L'énoncé Lean est accepté par le noyau, avec les quatre contrôles de la porte ci-dessus & Que l'énoncé Lean traduise fidèlement l'énoncé mathématique ou physique visé \\
\tierL{} & Le résultat figure dans une source lue, référencée précisément (par exemple un théorème d'article publié) & Qu'il ait été reverifié par nos soins \\
\tierC{} & Un argument humain complet, écrit dans ce livre, relu par un agent d'IA indépendant & Toute vérification par machine, et toute relecture humaine \\
\tierX{} & Un calcul en virgule flottante, avec contrôles positif et négatif & Rien : une exploration numérique ne soutient jamais une affirmation \\
\bottomrule
\end{tabular}
\end{center}

Deux règles de composition, utilisées partout dans la suite.
\begin{enumerate}
\item \emph{Le niveau d'une conséquence est le plus faible des niveaux de ses hypothèses.} Un corollaire de \texttt{openai/math} est \tierA{} \emph{conditionnel} : le noyau a
vérifié l'implication, pas l'hypothèse amont, dont la preuve n'a pas été relue par des pairs.
\item \emph{Un résultat négatif est un résultat.} Quand une extension plausible échoue, on l'écrit avec son niveau.
\end{enumerate}

\begin{warnbox}[: l'adéquation de l'énoncé]
Le noyau certifie que l'énoncé \lean{universal_energy_minimum}, \emph{avec les définitions du dépôt}, découle des axiomes de la liste blanche. Il ne certifie
pas que \lean{DensityOne} soit la densité que le physicien a en tête, ni que \lean{energy} (une limite inférieure sur des disques fermés centrés) soit l'énergie de CKMRV. Aucun humain
n'a audité cette adéquation. C'est pourquoi chaque chapitre répète le niveau, et c'est la raison de la mention \og conditionnel \fg{}.
\end{warnbox}

\section*{Exercices}
\addcontentsline{toc}{section}{Exercices}
\label{ch1:sec:exos}

Les indications suivent la liste des énoncés. Les exercices sont de difficulté croissante (de * à ***).

\begin{exercise}[*]\label{ch1:ex:covol}
Montrer que la matrice de passage entre deux bases d'un même réseau est dans $\mathrm{GL}_2(\Z)$ et en déduire que le covolume ne dépend pas de la base.
Calculer le covolume et la distance minimale du réseau $\Z u\oplus\Z v$ avec $u=(1,0)$, $v=(1/2,\sqrt3/2)$, puis la normaliser à la densité $1$. Retrouver $a^2=2/\sqrt3$.
\end{exercise}

\begin{exercise}[*]\label{ch1:ex:voisins}
Dans $A$, déterminer tous les vecteurs de norme $a\sqrt3$ et leur nombre (indication : $m^2+mn+n^2=3$). Comparer au cas de $\Z^2$ (couche $|x|^2=2$). Justifier que $a>1$ : à covolume égal, la distance minimale de $A$ dépasse celle de $\Z^2$.
\end{exercise}

\begin{exercise}[**]\label{ch1:ex:comptage}
Écrire la démonstration de la proposition~\ref{ch1:prop:comptage} pour $\Lambda=\Z^2$ et en déduire que $\#\{(m,n)\in\Z^2 : m^2+n^2\le R^2\}=\pi R^2+O(R)$.
Cette erreur $O(R)$ est-elle optimale ? (Chercher la littérature sur le \emph{problème du cercle de Gauss}.)
\end{exercise}

\begin{exercise}[**]\label{ch1:ex:epstein}
Démontrer que $\sum_{x\in A\setminus0}|x|^{-s}$ diverge pour $s\le2$ et converge pour $s>2$. En utilisant $t^{-s/2}=\frac{\pi^{s/2}}{\Gamma(s/2)}\int_0^\infty\alpha^{s/2-1}e^{-\pi\alpha t}d\alpha$,
exprimer cette somme comme une intégrale de $\Theta_A(\alpha)-1$ contre $\alpha^{s/2-1}d\alpha$ et retrouver la condition $s>2$ à partir du comportement de $\Theta_A$ en $0$ et en $+\infty$.
\end{exercise}

\begin{exercise}[**]\label{ch1:ex:poisson}
(a) Démontrer la relation $\vartheta(1/\alpha)=\sqrt\alpha\,\vartheta(\alpha)$ de l'exemple~\ref{ch1:ex:jacobi} en refaisant la démonstration unidimensionnelle.
(b) En déduire un algorithme pour calculer $\vartheta(\alpha)$ à $10^{-12}$ près en ne sommant que $|n|\le 6$ quel que soit $\alpha>0$. (c) Calculer
$\Theta_{\Z^2}(1)=\vartheta(1)^2$ et comparer à la valeur $1{,}1803\ldots$ du texte.
\end{exercise}

\begin{exercise}[*]\label{ch1:ex:clotures}
Montrer que si $f$ est CM et $a>0$, $b\ge0$, alors $t\mapsto f(at+b)$ est CM. En déduire que $t\mapsto(t+1)^{-s}e^{-a t}$ est CM pour $s,a\ge0$.
Montrer par un calcul que $t\mapsto t^{-1}e^{-t}$ l'est aussi sans utiliser le théorème de composition.
\end{exercise}

\begin{exercise}[**]\label{ch1:ex:comp}
(a) Vérifier que $h(t)=\ln(1+t)$ est de Bernstein. (b) En déduire que $t\mapsto(\ln(1+t))^{-s}$ est CM pour $s\ge0$ (attention à la positivité de $h$ sur $]0,\infty[$).
(c) Montrer que $t\mapsto1/(1+\sqrt t)$ est CM en appliquant le corollaire~\ref{ch1:cor:sqrt}.
(d) Pourquoi l'hypothèse $h>0$ ne peut-elle pas être omise dans le théorème~\ref{ch1:thm:comp} ? (Considérer $h\equiv0$.)
\end{exercise}

\begin{exercise}[***]\label{ch1:ex:melange}
Soit $g(t)=t^{-1/2}-2(t+1)^{-1/2}+(t+4)^{-1/2}$. En remarquant que $g$ est un mélange de gaussiennes à \emph{poids signé}, expliquer pourquoi la positivité du poids échoue,
et retrouver par ce biais que $g$ n'est pas CM (comparer à $g(3)<0$).
\end{exercise}

\begin{exercise}[*]\label{ch1:ex:lean}
Parmi les quatre sorties suivantes de \lean{#print axioms}, lesquelles satisfont la porte d'axiomes ? Lesquelles satisfont la porte \emph{complète} à quatre contrôles ?
(i) \texttt{[propext, Classical.choice, Quot.sound]} ; (ii) \texttt{[propext, sorryAx]} ; (iii) \texttt{does not depend on any axioms} ; (iv) \texttt{[cheat, propext]}.
Expliquer pourquoi (iii), seule, ne suffit pas à conclure.
\end{exercise}

\begin{exercise}[**]\label{ch1:ex:fatou}
Soit $\Lambda$ un réseau et $g\ge0$ un potentiel quelconque. Montrer que l'énergie de disque $E_g(\Lambda;R)$ définie plus haut converge, quand $R\to\infty$,
vers $\sum_{a\in\Lambda\setminus0}g(|a|^2)$ (finie ou non). Ainsi, pour un réseau, énergie inférieure et énergie du réseau coïncident.
\end{exercise}

\subsection*{Indications}

\begin{itemize}
\item[\ref{ch1:ex:covol}.] Si $M,M'$ sont deux matrices de bases, $M'=MP$ avec $P$ entière et $M=M'P'$ avec $P'$ entière : $PP'=I$, donc $\det P\det P'=1$ avec deux
entiers, donc $\det P=\pm1$. Pour la fin : $\det(u,v)=\sqrt3/2$ ; on met à l'échelle de rapport $a$ avec $a^2\sqrt3/2=1$.
\item[\ref{ch1:ex:voisins}.] $m^2+mn+n^2=3$ : $(1,1)$, $(-1,-1)$, $(2,-1)$, $(-2,1)$, $(1,-2)$, $(-1,2)$, soit six vecteurs. Comme $a^2=2/\sqrt3>1$ ($4>3$), $a>1$ ; la
distance minimale de $\Z^2$ est $1$.\item[\ref{ch1:ex:comptage}.] Prendre $\delta=\sqrt2$ pour $\Z^2$ : $\pi(R-\sqrt2)^2\le N(R)\le\pi(R+\sqrt2)^2$, d'où $N(R)=\pi R^2+O(R)$ ; l'optimalité de l'exposant $1$ n'est pas demandée (le meilleur exposant connu est inférieur à $1$, voir la littérature).
\item[\ref{ch1:ex:epstein}.] Pour la divergence, reprendre la minoration de $\#S_k$. Pour l'intégrale : $\Theta_A(\alpha)-1\sim6e^{-\pi a^2\alpha}$ quand $\alpha\to\infty$ (convergence)
et $\Theta_A(\alpha)\sim\alpha^{-1}$ quand $\alpha\to0$ (par Poisson) ; l'intégrale $\int_0\alpha^{s/2-1}\alpha^{-1}d\alpha$ converge ssi $s/2-2>-1$, i.e. $s>2$.
\item[\ref{ch1:ex:poisson}.] (b) Pour $\alpha\ge1$ on somme directement ; pour $\alpha<1$ on applique $\vartheta(\alpha)=\alpha^{-1/2}\vartheta(1/\alpha)$.
(c) $\vartheta(1)=1{,}0864348\ldots$, donc $\Theta_{\Z^2}(1)=1{,}18034\ldots$
\item[\ref{ch1:ex:clotures}.] $f(at+b)$ : dériver $n$ fois et factoriser $a^n$. Pour la dernière : produit de deux fonctions CM (proposition~\ref{ch1:prop:cloture}) ; ou calcul direct par Leibniz.
\item[\ref{ch1:ex:comp}.] (b) Poser $f(u)=u^{-s}$, CM. (d) Pour $h\equiv 0$, $f\circ h$ n'est pas défini sur $]0,\infty[$ (ou vaut une valeur parasite $f(0)$ si on prolonge).
\item[\ref{ch1:ex:melange}.] $t^{-1/2}=\pi^{-1/2}\int\alpha^{-1/2}e^{-\alpha t}d\alpha$ ; le poids de $g$ est $\pi^{-1/2}\alpha^{-1/2}(1-2e^{-\alpha}+e^{-4\alpha})$, qui est négatif pour
certains $\alpha$ (étudier $1-2x+x^4$ avec $x=e^{-\alpha}$). L'unicité de $\mu$ dans Bernstein--Widder et la positivité de $g$ en $t=3$ concluent.
\item[\ref{ch1:ex:lean}.] (i) passe. (ii) échoue (\lean{sorryAx}). (iii) passe la porte d'axiomes seule, mais la déclaration pourrait avoir été insérée en contournant le
vérificateur de types du noyau. (iv) échoue.
\item[\ref{ch1:ex:fatou}.] Le terme de $x\in\Lambda_R$ est $\sum_{a\in(\Lambda_R-x)\setminus0}g(|a|^2)$ : il est majoré par $\sum_{a\ne0}g(|a|^2)$, d'où la majoration de $E_g(\Lambda;R)$.
Pour la minoration, fixer $r>0$ : si $|x|\le R-r$, le terme de $x$ contient tous les $a\ne0$ avec $|a|\le r$. Donc
$E_g(\Lambda;R)\ge\frac{N_\Lambda(R-r)}{N_\Lambda(R)}\sum_{0<|a|\le r}g(|a|^2)$ ; le rapport tend vers $1$ (proposition~\ref{ch1:prop:comptage}). Faire tendre $R$ puis $r$ vers l'infini.
\end{itemize}
